\chapter{Projeto}
\label{capitulo:projeto}

Este capítulo apresenta os principais elementos de projeto do sistema proposto.
Neste sentido são definidas estruturas da arquitetura do sistema, protótipos e diagramas estruturais e comportamentais.

Embora este não seja o momento de definir conceitos\footnote{O momento de definir conceitos importantes é na revisão bibliográfica, capítulo \ref{capitulo:revisao_bibliografica}.}, você pode endossar o porquê as coisas foram feitas do jeito que foram utilizando citações.

\section{Arquitetura do projeto}

Fazer um diagrama de implantação da UML ou fazer uma figura que expresse como as partes do sistema foram divididas e como elas se comunicam.

\section{Protótipos}
\label{secao:prototipos}

Nesta seção são apresentados os principais, senão todos, os protótipos do sistema.
Conceitue o que são protótipos e lembre que eles são rascunhos do sistema, não as telas finais.
As telas finais são tema do capítulo de resultados e discussão (capítulo \ref{capitulo:resultados_discussao}).

Apresente todos os protótipos sempre com uma explicação e a figura.

\section{Diagrama Entidade Relacionamento (DER)}

Nesta seção encontra-se um exemplo de conceito que normalmente não entra no capítulo de revisão, mas cai muito bem uma boa citação.

Por que é importante citar aqui? 
Veja qual o motivo de ter escolhido um banco de dados relacional? 
E por que a modelagem foi feita da forma que foi? 
Ela segue as normas formais? 
Quantas normas formais são? 
Segue todas elas?

Então, além de apresentar a figura do DER você também deve fazer uma explicação e afirmar que ela segue as normas formais. Em casos que não seguirem as normas formais deve ser prontamente justificado o porquê.

É importantíssimo conferir a coesão do DER com os protótipos do sistema apresentados na seção \ref{secao:prototipos}.
Todos os campos apresentados no protótipo devem estar no DER.

\section{Diagrama de Casos de Uso}

Inclua e explique o diagrama de caso de uso.

\section{Diagrama de Classes}

Inclua e explique o diagrama de caso de classes.
Note que o diagrama de classes pode ter classes de domínio e sua relação assim como também podem haver diagramas de classes com mais níveis de detalhamento.
Caso seja necessário inclua mais detalhes.

\section{Algoritmos propostos}

No seu trabalho você não deve incluir códigos\footnote{Raramente isto é aceitável}.
Dê preferência a incluir os algoritmos que representam contribuições claras do seu trabalho.
Veja se está fazendo um sistema de \textit{e-commerce}, inclua o algoritmo de indicação de produtos, por exemplo o algoritmo \ref{alg:indicacao_produtos}. E se o sistema for de biblioteca, inclua o algoritmo de sugestão de livros.
E assim vai, converse com seu orientador o que seriam os algoritmos mais relevantes do seu trabalho.

\begin{algorithm}
\caption{Algoritmo de indicação de produtos.}
\label{alg:indicacao_produtos}

$P$ \gets $pesquisar\_produtos()$\;

$n \gets length(produtos)\;

\eSe{$n \geq 1$}
{
    $P_{ord}$ \gets $ordenar\_por\_estrelas(P)$\;
    
    $P_{preferidos}$ \gets $P_{ord}[1:10]$\;

    \Retorna{$P_{preferidos}$}
}{
    \Retorna{Lista vazia}
}
\end{algorithm}